\section{Maximal Roe algebras and the coarse index}
\label{sec:maximal-index}

The infinite-order class is now constructed. It remains to show that its
two end contributions have the same maximal index. The subsequent averaging
and sliding arguments will use operator
formulas on the standard modules.  We therefore establish the norm estimates
that allow those formulas to be used in the maximal completion, and then
verify the required naturality of the index.  These are standard features of
maximal Roe theory; see \cite[Section~2]{OOY09} and
\cite[Section~2]{SS21}.  We give the details needed here because the distinction
between operator norm and maximal norm is essential to the proof.

\subsection{Norm control and multipliers}

Let $Z$ be a uniformly discrete proper metric space of bounded geometry, and
write
\[
 H_0=\ell^2(\N),\qquad H_Z=\ell^2(Z)\otimes H_0,
 \qquad A_Z=C^*_{\max}(Z).
\]
The algebra $\C[Z]$ consists of the bounded finite-propagation operators on
$H_Z$ that are locally compact for the multiplication representation of
$C_0(Z)$.  We denote their concrete operator norm by $\|\cdot\|_{\mathrm{op}}$
and their maximal norm by $\|\cdot\|_{\max}$.  Write $T_{xy}:H_0\to H_0$ for
the matrix entries of an operator $T$.

For a bounded finite-propagation operator, local compactness is equivalent
to compactness of every entry $T_{xy}$.  In one direction, compress by the
characteristic functions of single vertices.  Conversely, every row and
column of a finite-propagation matrix has only finitely many entries.
Multiplication on either side by a finitely supported function therefore
gives a finite matrix of compact operators.  Approximation by finitely
supported functions proves the assertion for every function in $C_0(Z)$.
In particular, no uniform bound on the ranks, or on finite-rank
approximations to the entries, is part of this definition.

\begin{lemma}[Fixed-propagation norm control]
\label{ana:fixed-propagation}
For every $R\geq 0$ there is a constant $c_Z(R)<\infty$ such that
\[
 \|T\|_{\mathrm{op}}\leq\|T\|_{\max}
 \leq c_Z(R)\|T\|_{\mathrm{op}}
 \qquad(T\in\C[Z],\ \prop(T)\leq R).
\]
Consequently, operator-norm convergence of algebraic Roe operators with one
common propagation bound implies convergence in maximal norm.
\end{lemma}

\begin{proof}
Put $N_R=\sup_{z\in Z}|B_Z(z,R)|$.  The support of $T$, viewed as a bipartite
graph with a row copy and a column copy of $Z$, has degree at most $N_R$.
A greedy edge coloring partitions it into at most $2N_R-1$ matchings.
Accordingly, $T=\sum_{a=1}^{2N_R-1}T^{(a)}$, where each $T^{(a)}$ has at most
one nonzero entry in each row and column, and
$\|T^{(a)}\|_{\mathrm{op}}\leq\|T\|_{\mathrm{op}}$.

The diagonal algebra
\[
 \prod_{z\in Z}\mathcal K(H_0)
 =\left\{\operatorname{diag}(k_z): k_z\in\mathcal K(H_0),\
                  \sup_z\|k_z\|<\infty\right\}
\]
is a $C^*$-subalgebra of $\C[Z]$.  Every $*$-representation of $\C[Z]$
restricts to a contractive representation of this algebra.  Since
$(T^{(a)})^*T^{(a)}$ is diagonal, every such representation $\pi$ satisfies
\[
 \|\pi(T^{(a)})\|^2
 \leq\|(T^{(a)})^*T^{(a)}\|_{\mathrm{op}}
 \leq\|T\|_{\mathrm{op}}^2.
\]
The upper bound follows with $c_Z(R)=2N_R-1$.  The concrete representation
on $H_Z$ gives the lower bound.  Apply the estimate to differences to
obtain the convergence assertion.
\end{proof}

This matching argument is the discrete-module form of the maximal-norm
estimate in \cite[Lemma~2.1]{OOY09}; compare
\cite[Lemma~1.10]{SW13}.  We shall also use the following stronger statement
when the vertex support, rather than merely the propagation, is fixed.

\begin{lemma}[The compact ideal]
\label{ana:compacts}
An operator in $\C[Z]$ supported on $F\times F$, for a finite set $F\subset Z$,
has the same maximal and concrete norms.  The closure in $A_Z$ of all such
operators is an ideal $\mathcal K_Z$, canonically isomorphic to
$\mathcal K(H_Z)$.  In particular, $K_1(\mathcal K_Z)=0$.
\end{lemma}

\begin{proof}
The operators supported on $F\times F$ form the $C^*$-algebra
$\mathcal K(\ell^2(F)\otimes H_0)$.  Restriction of any representation is
contractive, and the concrete representation gives the reverse inequality.
These finite-vertex algebras have dense union in $\mathcal K(H_Z)$.
Multiplying one of their elements on either side by a finite-propagation
operator with compact entries gives another compact operator of finite
vertex support: only a finite metric neighborhood of $F$ can occur.
Density proves the ideal assertion.
\end{proof}

We will need finite-propagation operators whose entries are bounded but
need not be compact, such as the unitaries implementing the group action.
Here and below $M(A_Z)$ denotes the multiplier algebra of $A_Z$.

\begin{lemma}[Finite-propagation multipliers]
\label{ana:multipliers}
Every bounded finite-propagation operator $D$ on $H_Z$ acts by left and
right multiplication on $\C[Z]$, and these actions extend to a multiplier
of $A_Z$.  This assignment is a faithful $*$-homomorphism and satisfies
\[
 \|D\|_{\mathrm{op}}\leq \|D\|_{M(A_Z)}
 \leq c_Z(R)\|D\|_{\mathrm{op}}
 \qquad\text{if }\prop(D)\leq R.
\]
For an operator supported on a single matching, the multiplier norm equals
the concrete norm.
\end{lemma}

\begin{proof}
Products of $D$ with elements of $\C[Z]$ have finite propagation and compact
entries.  Decompose $D$ into matchings as in
Lemma~\ref{ana:fixed-propagation}.  It suffices to show that a matching
operator whose entries have norm at most one acts contractively.
For such an operator, $D^*D$ is a diagonal positive contraction.  The
operator $W=(1-D^*D)^{1/2}$ is bounded and diagonal.  For $T\in\C[Z]$,
all terms in
\[
 T^*T-(DT)^*(DT)=(WT)^*(WT)
\]
belong to $\C[Z]$.  Applying each $*$-representation gives
$\|DT\|_{\max}\leq\|T\|_{\max}$.  Applying this to $D^*$ and taking
adjoints gives the corresponding right-multiplication bound.  The
algebraic compatibility identities between these left and right actions
pass to the completion, so they define a multiplier.  Summing the
matching estimates gives the asserted upper bound.

The assignment preserves products and adjoints on the dense algebra
$\C[Z]$.  The concrete representation of $A_Z$ on $H_Z$ is nondegenerate,
since it contains $\mathcal K(H_Z)$, and hence extends to the multiplier
algebra.  It sends the multiplier just constructed to the original
operator $D$, proving the lower bound and faithfulness.
\end{proof}

In particular, an operator identity involving bounded finite-propagation
operators remains an identity between their multipliers.  This statement
does not assert that the concrete representation of the entire maximal
Roe algebra is faithful.  The finite-propagation control in the preceding
lemmas is what will justify the specific limits below.

\subsection{Maps induced by coarse maps}

A map $h:Z\to W$ between uniformly discrete proper bounded-geometry spaces
is coarse if inverse images of bounded sets are bounded and, for every
$R$, there is $S$ such that $d_Z(x,y)\leq R$ implies
$d_W(h(x),h(y))\leq S$.  Choose an isometry
$V_h:H_Z\to H_W$ that sends the summand over $z$ into the summand over
$h(z)$, using orthogonal ranges when several vertices have the same image.
Such a choice is possible because the vertex fibers are finite and $H_0$
is infinite dimensional.  We call $V_h$ an exact covering isometry for $h$.

\begin{proposition}[Coarse functoriality]
\label{ana:functoriality}
Conjugation by an exact covering isometry defines a $*$-homomorphism
\[
 h_A=\Ad(V_h):A_Z\longrightarrow A_W.
\]
Its induced map on $K$-theory is independent of the covering isometry.
These induced maps respect composition, and an exact covering isometry
over $\id_Z$ induces the identity on $K_*(A_Z)$.
\end{proposition}

\begin{proof}
For $T\in\C[Z]$, coarse control bounds the propagation of $V_hTV_h^*$.
Each of its vertex entries is a finite sum of compact operators, because
$h^{-1}(w)$ is finite for each $w$.  Thus conjugation is a $*$-homomorphism
$\C[Z]\to\C[W]$, and the universal property extends it contractively to
the maximal completions.

For two choices $V_0,V_1$, put
\[
 V_\lambda=
 \begin{pmatrix}\cos\lambda\,V_0\\ \sin\lambda\,V_1\end{pmatrix}
 :H_Z\longrightarrow H_W\oplus H_W,
 \qquad 0\leq\lambda\leq\pi/2.
\]
Each $V_\lambda$ is an isometry.  Conjugation defines a homomorphism to
$M_2(A_W)$: its four matrix entries have compact vertex entries and one
common propagation bound for every fixed $T\in\C[Z]$.  Moreover,
\[
 \|V_\lambda T V_\lambda^*-V_\nu T V_\nu^*\|_{\mathrm{op}}
 \leq 2\|T\|_{\mathrm{op}}|\lambda-\nu|.
\]
Lemma~\ref{ana:fixed-propagation}, applied to the four entries, promotes
this estimate to maximal-norm continuity on $\C[Z]$; contractivity and
density give continuity on $A_Z$.  The two endpoint homomorphisms are
$\Ad(V_0)$ and $\Ad(V_1)$ in the two standard matrix corners, which induce
the same stabilized $K$-theory map.  Composing covering isometries proves
functoriality, and comparison with the identity isometry proves the last
assertion.
\end{proof}

This is the exact-covering version of \cite[Lemma~2.15]{OOY09}.  For the
maps in our construction we henceforth choose such homomorphisms and write
\[
 f_A:A_Y\longrightarrow A_B,\qquad
 i_{\pm,A}:A_B\longrightarrow A_X.
\]
By Proposition~\ref{ana:functoriality}, all statements about their
$K$-theory maps are independent of these choices.

\subsection{The canonical odd index and naturality}

The remaining index-theoretic input is compatibility with the two end
inclusions.

\begin{proposition}[Naturality for the end classes]
\label{ana:naturality}
For the classes $\beta\in\KX_1(B)$ and
$\alpha=(i_+)_*\beta-(i_-)_*\beta\in\KX_1(X)$ constructed above,
\begin{equation}
 \label{eq:ana-naturality}
 \mu_X^{\max}(\alpha)
 =(i_{+,A})_*\mu_B^{\max}(\beta)
  -(i_{-,A})_*\mu_B^{\max}(\beta).
\end{equation}
\end{proposition}

To verify this identity in the maximal completion, we first identify the
Roe algebra of a Rips complex with the vertex algebra and then recall the
exponential formula for its canonical odd index.

For $r>0$, regard a point of $P_r(Z)$ as a finitely supported probability
measure on $Z$ whose support has diameter at most $r$.  Equip the complex
with the $1$-Wasserstein metric induced by $d_Z$.  Every point is within
$r$ of every vertex of its support.  Moreover, a bounded set in this
metric uses only vertices from one bounded subset of $Z$: if $v$ is a
support vertex of $p$, then
\[
 d_Z(v,z_0)\leq r+\int_Z d_Z(z,z_0)\,dp(z)
             =r+d_{W_1}(p,\delta_{z_0}).
\]
That subset is finite.  Consequently this metric is proper and induces
the usual topology on the locally finite complex.  The vertex inclusion
and any selection of a support vertex are coarse inverses.

Choose a countable dense subset $E\subset P_r(Z)$ containing the vertices,
and use the ample module $\ell^2(E)\otimes H_0$ with pointwise
multiplication by $C_0(P_r(Z))$.  Assign each point of $E$ to one vertex of
its support, assigning a vertex to itself.  This partitions the module
into infinite-dimensional separable blocks $H_v$, one for each $v\in Z$;
choose identifications $H_v\cong H_0$.

Under this regrouping, the algebraic Roe algebra of $P_r(Z)$ is exactly
$\C[Z]$.  Propagation bounds in the two descriptions differ by at most
$2r$.  For local compactness, the points in $H_v$ lie in the compact
closed star of $v$, so a compactly supported continuous cutoff equal to
one on that star shows that a locally compact operator has compact
vertex-block entries.  Conversely, a compact subset of $P_r(Z)$ meets
only finitely many blocks.  Compact vertex-block entries and finite
propagation therefore give local compactness after multiplication by a
compactly supported function, and then by any function in $C_0(P_r(Z))$.
The algebraic identification extends to the maximal completions by
universality.  It is the coarse identification induced by the chosen
support-vertex map.  Pseudolocality below is always taken relative to
$C_0(P_r(Z))$; the regrouping is used only for the Roe algebra and for
Lemma~\ref{ana:multipliers}.

We recall the cycle description of the odd index.  On an ample module
for $P_r(Z)$, an odd $K$-homology class is represented, after stabilization
and addition of degenerate cycles, by a bounded self-adjoint operator
$F$ such that
\[
 [F,a]\text{ is compact},\qquad a(F^2-1)\text{ is compact}
 \quad(a\in C_0(P_r(Z))).
\]
The first condition is pseudolocality.  The use of a fixed ample module
is the usual Paschke realization of locally finite $K$-homology; see
\cite{HigsonRoe00}.

Such a representative may be taken to have finite propagation.  Indeed,
let $\varphi_v$ be the square roots of the barycentric coordinate
functions.  Their supports are compact, have diameter at most $2r$,
and form a locally finite family with $\sum_v\varphi_v^2=1$.  Set
\[
 F'=\sum_v\varphi_vF\varphi_v.
\]
The sum converges strongly and has norm at most $\|F\|$: it is the
compression of the direct sum of copies of $F$ by the isometry
$\xi\mapsto(\varphi_v\xi)_v$.  It has propagation at most $2r$.
Furthermore,
\[
 F'-F=\sum_v\varphi_v[F,\varphi_v]
\]
is locally compact, since multiplication on the left by a compactly
supported function leaves only finitely many compact summands; use
self-adjointness for the right-hand assertion.  Thus the straight-line
path from $F$ to $F'$ consists of cycles and does not change the
$K$-homology class.  Replace $F$ by $F'$.

After regrouping into vertex blocks, Lemma~\ref{ana:multipliers} makes
$F$ a self-adjoint multiplier of $A_Z$, and finite propagation gives
$F^2-1\in\C[Z]$.  Hence
\[
 p=\frac{1+F}{2}\in M(A_Z),\qquad
 p^2-p\in\C[Z],\qquad
 u_F=\exp(2\pi i p)\in 1+A_Z.
\]
The odd maximal index is
\begin{equation}
 \label{eq:ana-odd-index}
 \operatorname{Ind}_{Z,\max}(F)=[u_F]\in K_1(A_Z).
\end{equation}
To relate this formula to the canonical index, let
$D^*_{\max}(P_r(Z))$ be the universal completion of the algebra of
finite-propagation pseudolocal operators on the Rips module.  It contains
$A_Z$ as an ideal, and Paschke duality identifies
\[
 \KK_1(C_0(P_r(Z)),\C)
 \cong K_0\bigl(D^*_{\max}(P_r(Z))/A_Z\bigr),
 \qquad [F]\longmapsto[p\bmod A_Z].
\]
These assertions, including existence of the universal completion, are
proved in \cite[Proposition~2.22, Propositions~2.27--2.28, and
Corollary~2.29]{SS21}.  The index is the boundary map of this extension.
The multiplier action of Lemma~\ref{ana:multipliers} agrees with the
inclusion on $A_Z$, so its exponential boundary formula is exactly
\eqref{eq:ana-odd-index}.  Thus this formula computes the canonical index,
with its usual boundary-map sign convention.  It agrees with the cycle
construction in \cite[Corollary~2.18]{OOY09}.  Passing to the Rips direct
limit gives $\mu_Z^{\max}$.



\begin{proof}[Proof of Proposition~\ref{ana:naturality}]
Both classes are represented at Rips scale $3$, where the end maps are
proper continuous coarse simplicial inclusions
$i_{\pm,3}:P_3(B)\to P_3(X)$.  Choose the dense-set modules compatibly:
the target dense set contains the images of the source dense set, and
the support-vertex choices agree on these images.  The resulting module
isometry $V$ is an exact covering isometry for the end map after
regrouping.

If $F$ represents a source odd class and has finite propagation, its
push-forward is represented on the target module by
\[
 \widetilde F=VFV^*-(1-VV^*).
\]
The complement is a degenerate cycle.  Properness ensures that pullbacks
of $C_0$ functions lie in $C_0$, so $\widetilde F$ is pseudolocal and
$\widetilde F^2-1$ is locally compact.  Coarse control preserves finite
propagation.  Writing $p=(1+F)/2$, we have
$\widetilde p=(1+\widetilde F)/2=VpV^*$.

For completeness, compatibility of the exponential with conjugation
holds in the maximal completion itself.  The function
\[
 h(z)=\frac{\exp(2\pi i z)-1}{z^2-z}
\]
extends to an entire function.  Approximate $h$ uniformly by polynomials
on disks containing the spectra of the two multipliers.  Since
\[
 \exp(2\pi i p)-1=(p^2-p)h(p),
\]
and $p^2-p\in\C[B]$, all polynomial identities are identities on the
algebraic Roe algebras, and their limits hold in maximal norm.  They give
\[
 \exp(2\pi i\widetilde p)-1
   =i_{\pm,A}\bigl(\exp(2\pi i p)-1\bigr).
\]
Thus the two indices commute with the end inclusion.  Apply this to
$\beta_3$, subtract the two identities, and use the definition of the
Rips direct-limit assembly map.  Independence of the covering choices
follows from Proposition~\ref{ana:functoriality}.
\end{proof}
